\begin{appendices}
\addtocontents{toc}{\protect\setcounter{tocdepth}{1}}
\section{Calculation of curved spacetime expansions from flat spacetime expansions} \label{Horizons:app:ExpansionCalculation} 
In this appendix we give a detailed derivation of our main result relating the expansions on the $\eta_{\mu \nu}$ and the $g_{\mu \nu}$ spacetimes \eqref{thetaFromTheta0} for a generic Kerr-Schild metric
\eqref{kerrSchildMetric}
\begin{equation}
    g_{\mu \nu} = \eta_{\mu \nu} + \Phi k_{\mu} k_{\nu}.
\end{equation}
We will use coordinate-basis methods throughout. Computing expansions is also possible in an orthonormal basis using the Newman-Penrose formalism, where the expansion appears as one of the various spin coefficients. Our main result \eqref{thetaFromTheta0} is corroborated by just such a calculation of ``shifted spin coefficients'' performed by \cite{BenTov:2017kyf} in a different context.

First, for notational simplicity, we define the vectors $l^{(0)\mu}$ and $l^{\mu}$, which live respectively on the $\eta_{\mu \nu}$ and $g_{\mu \nu}$ backgrounds,
\begin{align}
    l^{(0) \mu} &\coloneqq \eta^{\mu \nu} l^{(0)}_{\nu} \\
    l^{\mu} &\coloneqq g^{\mu \nu} l_{\mu}.
\end{align}
Thus, including \eqref{lFroml0Down}, we have the relations
\begin{align}
    l^{\mu} &= l^{(0) \mu} + \frac{\Phi}{2} k^{\mu} \label{Horizons:lFroml0Up} \\
    l_{\mu} &= l^{(0)}_{\mu} - \frac{\Phi}{2} k_{\mu}.
\end{align}
In terms of $l^{\mu}$ and $l^{(0) \mu}$, we can write the induced (inverse) metrics $h^{\mu \nu}$ \eqref{flatInducedMetric} and $\mathring{h}^{\mu \nu}$ \eqref{curvedInducedMetric} more succinctly as
\begin{align}
    h^{\mu \nu} &= g^{\mu \nu} + k^\mu l^\nu + l^\mu k^\nu \\
    \mathring{h}^{\mu \nu} &= \eta^{\mu \nu} + k^\mu l^{(0) \nu} + l^{(0) \mu} k^\nu.
\end{align}
Using the explicit inverse Kerr-Schild metric \eqref{kerrSchildMetricInverse} and the relation \eqref{lFroml0Up}, 
\begin{equation}
    \begin{split}
        h^{\mu \nu}
        &= \eta^{\mu \nu} - \Phi k^\mu k^\nu + k^\mu \left( l^{(0) \nu} + \frac{\Phi}{2} k^\nu \right) + \left( l^{(0) \mu} + \frac{\Phi}{2} k^\mu \right) k^\nu \\
        &= \eta^{\mu \nu} + k^\mu l^{(0) \nu} + l^{(0) \mu} k^\nu \\
        &= \mathring{h}^{\mu \nu}.
    \end{split}
\end{equation}
We have thus found that the metric on $\mathcal{S}_\lambda$ induced from the flat and the curved background are equal. We henceforth drop the circular accent on $\mathring{h}^{\mu \nu}$.

To relate the flat- and curved-spacetime expansions defined in \eqref{flatExpansions} and \eqref{curvedExpansions}, we will also need to relate the derivatives $\mathring{\nabla}_\mu$ and $\nabla_\mu$. It is a standard result in differential geometry that the difference between the action of any two covariant derivatives on a one-form is a linear map (see e.g. chapter 3 of \cite{Wald:1984rg}), 
\begin{equation}
    \nabla_\mu V_\rho - \mathring{\nabla}_\mu V_\rho = -\Gamma_{\mu \rho}^{\sigma} V_{\sigma}.
\end{equation}
We can thus write
\begin{align}
    \begin{split} \label{Horizons:thetaKIntermediate}
        \theta^{(k)}
        &= h^{\mu \nu} \nabla_\mu k_\nu \\
        &= h^{\mu \nu} \mathring{\nabla}_\mu k_\mu - h^{\mu \nu} \Gamma_{\mu \nu}^\sigma k_\sigma \\
        &= \mathring{\theta}^{(k)} - h^{\mu \nu} \Gamma_{\mu \nu}^\sigma k_\sigma,
    \end{split} \\
    \intertext{and}
    \begin{split} \label{Horizons:thetaLIntermediate}
        \theta^{(l)}
        &= h^{\mu \nu} \nabla_\mu l_\nu \\
        &= h^{\mu \nu} \mathring{\nabla}_\mu l_\nu - h^{\mu \nu} \Gamma_{\mu \nu}^\sigma l_\sigma \\
        &= h^{\mu \nu} \mathring{\nabla}_\mu \left(l_\nu^{(0)} - \frac{\Phi}{2} k_\nu \right) - h^{\mu \nu} \Gamma_{\mu \nu}^\sigma l_\sigma \\
        &= \mathring{\theta}^{(l^{(0)})} -  h^{\mu \nu} \Gamma_{\mu \nu}^\sigma l_\sigma - h^{\mu \nu} \mathring{\nabla}_\mu \left( \frac{\Phi}{2} k_\sigma \right) \\
        &= \mathring{\theta}^{(l^{(0)})} -  h^{\mu \nu} \Gamma_{\mu \nu}^\sigma l_\sigma - \frac{\Phi}{2}\mathring{\theta}^{(k)}.
    \end{split}
\end{align}
We demonstrated that the last equality holds in section \ref{sec:doubleCopyRelation}. Clearly, we will need the contractions $h^{\mu \nu} \Gamma_{\sigma \mu \nu} k^\sigma$ and $h^{\mu \nu} \Gamma_{\sigma \mu \nu} l^\sigma$.

First, from the metric compatibility of the two derivative operators with $g_{\mu \nu}$ and $\eta_{\mu \nu}$ respectively, we can write
\begin{equation}
    \Gamma_{\sigma \mu \nu} = g_{\sigma \lambda} \Gamma^\lambda_{\mu \nu} = \frac{1}{2}\left( \mathring{\nabla}_{\nu} \left[\Phi k_{\mu} k_{\sigma}\right] +
    \mathring{\nabla}_{\mu} \left[\Phi k_{\nu} k_{\sigma}\right] - \mathring{\nabla}_{\sigma}\left[\Phi k_{\mu} k_{\nu}\right]\right).
\end{equation}
Let us simplify $h^{\mu \nu}\Gamma_{\sigma \mu \nu}$,
\begin{equation} \label{Horizons:hGamma}
    h^{\mu \nu} \Gamma_{\sigma \mu \nu} = \eta^{\mu \nu} \Gamma_{\sigma \mu \nu} + 2 k^\mu l^{(0)\nu} \Gamma_{\sigma \mu \nu},
\end{equation}
where we used the index symmetry of $\Gamma_{\sigma \mu \nu}$. Using the metric compatibility of $\mathring{\nabla}$ with $\eta_{\mu \nu}$ and the nullness of $k_\mu$, the first term in \eqref{hGamma} simplifies as
\begin{equation}
    \frac{1}{2} \eta^{\mu \nu}\left( \mathring{\nabla}_{\nu} \left[\Phi k_{\mu} k_{\sigma}\right] +
\mathring{\nabla}_{\mu} \left[\Phi k_{\nu} k_{\sigma}\right] - \mathring{\nabla}_{\sigma}\left[\Phi k_{\mu} k_{\nu}\right]\right) = \mathring{\nabla}_\mu \left[ \Phi k^\mu k_\sigma \right].
\end{equation}
To simplify the second term of \eqref{hGamma}, we compute
\begin{equation}  
    \begin{split}
        k^{\mu} \Gamma_{\sigma \mu \nu} &= \frac{1}{2} k^{\mu} \mathring{\nabla}_{\mu} \left[\Phi k_{\nu} k_{\sigma}\right] + \frac{1}{2} k^{\mu}
        \mathring{\nabla}_{\nu} \left[\Phi k_{\mu} k_{\sigma}\right] - \frac{1}{2} k^{\mu} \mathring{\nabla}_{\sigma} \left[\Phi k_{\mu} k_{\nu}\right] 
        \\
                           &= \frac{1}{2} \Phi k_{\nu} k^{\mu} \mathring{\nabla}_{\mu}k_{\sigma} + \frac{1}{2} \Phi k_{\sigma} k^{\mu}
                           \mathring{\nabla}_{\mu} k_{\nu} + \frac{1}{2} k_{\nu} k_{\sigma} k^{\mu} \mathring{\nabla}_{\mu} \Phi + \frac{1}{2}
                           k^{\mu}k_{\mu} \mathring{\nabla}_{\nu}\left[\Phi k_{\sigma}\right] \\
                           &\phantom{=\ }\ + \frac{1}{2} \Phi k_{\sigma} k^{\mu}
                           \mathring{\nabla}_{\nu} k_{\mu} - \frac{1}{2} k^{\mu} k_{\mu} \mathring{\nabla}_{\sigma} \left[\Phi k_{\nu}\right] -
                           \frac{1}{2} \Phi k_{\nu} k^{\mu} \mathring{\nabla}_{\sigma} k_{\mu} \\
                           &= \frac{1}{2} k_{\nu} k_{\sigma} k^{\mu} \mathring{\nabla}_{\mu}\Phi.
    \end{split}
\end{equation}
The last line here follows from the previous by observing that the nullness and geodeticity of $k_\mu$ forces all
but the third term to vanish. The second term of \eqref{hGamma} is therefore
\begin{equation}
    2 k^\mu l^{(0) \nu} \Gamma_{\sigma \mu \nu} = l^{(0) \nu}  k_\nu k_\sigma k^\mu \mathring{\nabla}_\mu \Phi = - k_\sigma k^\mu \mathring{\nabla}_\mu \Phi,
\end{equation}
where we used the normalization condition $l^{(0)\mu} k_\mu = -1$. We have thus simplified \eqref{hGamma} to
\begin{equation}
    \begin{split}
        h^{\mu \nu} \Gamma_{\sigma \mu \nu}
        &= \mathring{\nabla}_\mu \left[ \Phi k^\mu k_\sigma \right] - k_\sigma k^\mu \mathring{\nabla}_\mu \Phi \\
        &= \Phi \mathring{\nabla}_\mu \left[ k^\mu k_\sigma \right] \\
        &= \Phi k_\sigma \mathring{\nabla}_\mu k^\mu,
    \end{split}
\end{equation}
where to obtain the final equality we used the geodeticity of $k_\sigma$. In just a moment we will need also the equality $\mathring{\theta}^{(k)} = \mathring{\nabla}_\mu k^\mu$, which we now show. The difference between the two terms can be simplified as,
\begin{equation}
    \begin{split}
        \mathring{\theta}^{(k)} - \mathring{\nabla}_\mu k^\mu
        &= \left( h^{\mu \nu} - \eta^{\mu \nu} \right) \mathring{\nabla}_\mu k_\nu \\
        &= \left(k^\mu l^{(0)\nu} + l^{(0)\mu} k^\nu \right) \mathring{\nabla}_\mu k_\nu \\
        &= l^{(0) \nu} k^\mu \mathring{\nabla}_\mu k_\nu + \frac{1}{2} l^{(0) \mu} \mathring{\nabla}_\mu \left(k^\nu k_\nu \right) \\
        &= 0,
    \end{split}
\end{equation}
where the first term vanishes due to geodeticity and the second due to the nullness of $k^\mu$. We can finally compute the two contractions we needed,
\begin{align}
    h^{\mu \nu} \Gamma_{\sigma \mu \nu} k^\sigma &= \Phi \left(k^\sigma k_\sigma\right) \mathring{\nabla}_\mu k^\mu = 0, \label{Horizons:hGammaK} \\
    h^{\mu \nu} \Gamma_{\sigma \mu \nu} l^\sigma &= \Phi \left(l^\sigma k_\sigma\right) \mathring{\nabla}_\mu k^\mu = -\Phi \mathring{\theta}^{(k)}. \label{Horizons:hGammaL}
\end{align}
Substituting \eqref{hGammaK} into \eqref{thetaKIntermediate} and \eqref{hGammaL} into \eqref{thetaLIntermediate},
\begin{align}
    \theta^{(k)} &= \mathring{\theta}^{(k)} - h^{\mu \nu}\Gamma_{\sigma \mu \nu} k^\sigma = \mathring{\theta}^{(k)}. \\
    \begin{split}
        \theta^{(l)} &= \mathring{\theta}^{(l^{(0)})} -  h^{\mu \nu} \Gamma_{\mu \nu}^\sigma l_\sigma - \frac{\Phi}{2}\mathring{\theta}^{(k)} = \mathring{\theta}^{(l^{(0)})} + \Phi\mathring{\theta}^{(k)} - \frac{\Phi}{2}\mathring{\theta}^{(k)} \\
        &= \mathring{\theta}^{(l^{(0)})} + \frac{\Phi}{2}\mathring{\theta}^{(k)}.
    \end{split}
\end{align}

\end{appendices}

